\section{Related Work}

As discussed in previous studies~\cite{googleFinetuning2025}, spatial understanding for assistive technology is an active research area.

\subsection{Spatial Understanding Systems}

\subsubsection{SpatialBot}
SpatialBot is a VLM designed for spatial tasks, by incorporating depth images and data into its outputs.
It is trained on a hierarchy of spatial tasks, with are categorized as low, middle and high-level.
\begin{itemize}
    \item \textbf{Low-level tasks:} Direct querying of depth values from a depth map and exact information of these on set points
    \item \textbf{Middle-level tasks:} Comparing distances between objects, and calculating the relative position of objects in a scene.
    \item \textbf{High-level tasks:} Understanding spatial relationships between objects and doing things like counting objects in a scene.
\end{itemize}

It has been successfully deployed in robotic systems, where it can be used to translate spatial understanding into physical actions.
This highlights the robustness of its capabilities with would help in guidance systems, where it is important to get correct spatial information to the user.

Another feature that is highlighted in the paper is the ability to distinguish between real and printed objects by using the spatial depth information provided by the depth images.
This can also help in navigation where detecting hazards is important, with also means the system needs to be able to decide if an object is a poster or a real thread.

Internally the SpatialBot framework uses a fine-tuned VLM with a curated RGB-D dataset.
It tries to specifically address the challenge of providing consistent and accurate information for robotic systems.
As seen later in this paper, we also used a lot of the same concepts in mapping depth values and providing spatial information als seen here.\cite{SpatialBotCai2025}

\subsubsection{Video-3D LLM}
Video-3D LLM is a model that enhances Multimodal Large Language Models (MLLMs) with sophisticated 3D spatial understanding capabilities.
This tackles the challenge, that most MLLMs are only trained on 2D data with limits there capabilities in spatial reasoning.
It is treating 2D scenes as dynamic videos and incorporating 3D data into the video representation.
This allows a more accurate understanding and goes in line with the goal of using a smartphone video stream as an input for our system.

\begin{itemize}
    \item \textbf{Position-Aware Video Representation:} It is not using the traditional approach of converting 3D scenes into point clouds. Instead it directly processes video frames along with their corresponding 3D global coordinates. It does this by back-projecting pixel from the depth image and then encoding and injecting the visual features to give the option of getting spatial information.
    \item \textbf{Efficient Frame Sampling:} To handle larger video sequences, it uses a sampling strategy where it only uses the the minimum number of frames needed to get all the spatial information.
    \item \textbf{Multi-Task Generalist Model:} As a single model it can handle task like, describing objects, locating objects and answering questions about objects that require spatial reasoning.
\end{itemize}

Even though it might seem like a perfect fit for our usecase it still comes with some limitations.
For example, it is trained on indoor scenes and would probably not be able to handle outdoor scenes as well.
Its compute time for processing a video is also quite high, making it difficult to use in a real-time setting. \cite{zheng2024video3dllmlearningpositionaware}

\subsubsection{SpatialVLM}
\label{ssec:spatialvlm}

SpatialVLM~\cite{spatialvlmChen2024} by Google DeepMind is a data-synthesis and pre-training framework designed to endow generic vision-language models (VLMs) with \emph{quantitative} spatial reasoning capabilities.  It converts monocular RGB images into \emph{canonicalised}\footnote{The authors define \emph{canonicalisation} as a RANSAC-based alignment of the depth-derived point cloud to a physically meaningful world frame whose $z$-axis is normal to a detected horizontal support surface (e.g.\ floor or tabletop).} 3-D scene graphs and uses these graphs to auto-generate an internet-scale spatial VQA corpus.

\paragraph{System overview.}
The \emph{data synthesis pipeline} comprises four off-the-shelf expert models: CLIP to filter images of real world scenes, open-vocabulary detection, semantic segmentation, monocular metric depth estimation (ZoeDepth~\cite{ZoeDepthBhat2023}), and object-centric captioning.  Detected regions are lifted to a metric point cloud, canonicalised as described above, and clustered into object instances whose centroids, sizes, and pairwise Euclidean distances are computed, resulting in a spatial graph where objects are represented by vertices and relations between them are quantitative relationships. A templated question-answer generator then produces \(\sim\!2\text{billion}\) QA pairs over \(10\text{ million}\) real-world images, covering \textbf{38} qualitative and quantitative spatial relation types (e.g.\ \emph{left-of}, \emph{higher-than}, exact metric distance).

\paragraph{Generated dataset.}
The resulting VQA dataset mixes 50 \% qualitative and 50 \% quantitative questions.  \emph{One-sentence overview:} it spans fine-grained prompts such as ``\textit{How far is the blue mug behind the red plate?}'' or ``\textit{Which object is wider, the laptop or the book?}'', always answered in human-like units and rounding.

\paragraph{Downstream use.}
SpatialVLM's synthetic QA pairs were used to \emph{pre-train Googles Gemini models}, endowing them with foundational spatial reasoning capabilities that were evidentally extended to endow Gemini models with AABB and OBB detection, as well as semantic segmentation capabilities.

\paragraph{Criticism.}
While the framework captures metric distances and relative size/position, it omits \emph{relative orientations} (e.g.\ object yaw or facing direction).  Incorporating such labels—even if users rarely ask explicitly—could further enrich the learned inductive biases and enable reasoning about tasks like ``\emph{Is the chair facing the table?}'' or ``\emph{What's the yaw between the red car and the parked scooter?}''.

\subsection{Key Concepts and Limitations}
% Content about relevant concepts and limitations from existing work
% - Common approaches in spatial understanding
% - Identified gaps in current solutions
% - Challenges in real-world deployment
% - Opportunities for improvement